\begin{proofbox}
	\textbf{\textcolor{CProof}{思路提示}}
	二次函数一般式通过配方可得顶点式，再利用求根公式可分解为两根式，并由此得到根与系数的关系。

	\medskip
	\textbf{\textcolor{CProof}{正式推导}}
	\begin{description}[style=nextline, leftmargin=2.8em, labelsep=0.5em, itemsep=0.55em]
		\item[\textbf{步骤一（配方变形）：}]
			\[
				\begin{aligned}
					f(x) &= a\left(x^2+\frac{b}{a}x\right)+c \\ &= a\left(x+\frac{b}{2a}\right)^2 + c - \frac{b^2}{4a}.
			\end{aligned}
			\]
			\par\textbf{理由：} 这是配方法的标准步骤，在括号内加上并减去 $\frac{b^2}{4a^2}$ 再乘以 $a$，整理后得到顶点式。
		\item[\textbf{步骤二（顶点坐标）：}]
			\[
				f(x)=a(x-h)^2+k,\quad h=-\frac{b}{2a},\; k=\frac{4ac-b^2}{4a}.
			\]
			\par\textbf{理由：} 对比上一步，$h=-\frac{b}{2a}$，$k=c-\frac{b^2}{4a}$，通分即得 $k$ 的表达式。
		\item[\textbf{步骤三（求根公式）：}]
			当 $\Delta=b^2-4ac\ge0$ 时，求解方程 $ax^2+bx+c=0$ 得到两根
			\[
				x_{1,2}=\frac{-b\pm\sqrt{\Delta}}{2a}.
			\]
			\par\textbf{理由：} 直接应用一元二次方程的求根公式。
		\item[\textbf{步骤四（两根式分解）：}]
			根据因式定理，将两根代入分解式
			\[
				f(x)=a(x-x_1)(x-x_2).
			\]
			\par\textbf{理由：} 若 $x_1,x_2$ 是 $f(x)=0$ 的根，则 $f(x)$ 可被 $(x-x_1)(x-x_2)$ 整除，且二次项系数为 $a$。
		\item[\textbf{步骤五（和积关系）：}]
			展开两根式并与一般式对比系数
			\[
				(x-x_1)(x-x_2)=x^2-(x_1+x_2)x+x_1x_2,
			\]
			于是得
			\[
				x_1+x_2=-\frac{b}{a},\quad x_1x_2=\frac{c}{a}.
			\]
			\par\textbf{理由：} 展开后与 $ax^2+bx+c$ 除以 $a$ 后的二次项系数、一次项系数、常数项对应相等。
	\end{description}

	\medskip
	\textbf{\textcolor{CProof}{结论回扣}}
	以上推导证实了二次函数的一般式、顶点式、两根式三种形式等价，并给出具体的坐标公式与根与系数关系，解题时可灵活选用。
\end{proofbox}
